\providecommand{\Cscen}[1]{\mathcal C_{#1}}
\providecommand{\supp}{\operatorname{supp}}

\section{Sources shared by three or more observers}\label{sec:hypergraph}

Remark~\ref{rem:hypergraph} shows that the adjacency graph of the observers does not decide termination once a source may be read by three observers. This section gives the replacement. Termination is proved for every scenario that can be brought to a disjoint union of hyper-double-stars by deleting redundant sources and observers that read every source, and nontermination for every scenario containing a non-double-star pair-source scenario as a trace. On at most four observers the two results leave exactly one scenario undecided.

\subsection{Scenarios, reductions and traces}

\begin{definition}[Hypergraph scenario]\label{def:hscenario}
A \emph{hypergraph scenario} $H=(V,E)$ consists of a finite set $V$ of observers and a finite set $E$ of source labels, each label $e$ carrying a nonempty vertex set, also written $e\subseteq V$, such that every observer lies in some source. Distinct labels may carry the same vertex set. The scenario has one latent source $S_e$ for each $e\in E$, taking values in an arbitrary measurable space, and one observed variable $O_v$ with finite alphabet $\mathcal O_v$ for each $v\in V$; the sources are mutually independent and
\[
O_v\sim\kappa_v\bigl(\cdot\mid(S_e)_{e\ni v}\bigr),
\]
conditionally independently across $v$ given the sources. $\Cscen{H}$ is the set of joint laws of $(O_v)_{v\in V}$ so obtained. Write $E(v)=\{e\in E:v\in e\}$, $E(W)=\bigcup_{v\in W}E(v)$ and $\deg v=|E(v)|$.
\end{definition}

A pair-source scenario in the sense of Definition~\ref{def:pairsource} is the case in which every source has two members and distinct labels carry distinct vertex sets. The covering condition replaces the exclusion of isolated vertices: a response seed at $v$ can be appended to any source at $v$ and ignored by its other members, so deterministic and stochastic responses give the same $\Cscen{H}$.

The order-$t$ inflation and the sets $\INW_t(H)$, $\IAI_t(H)$ and $\Iexp_t(H)$ are defined by Definitions~\ref{def:ginflation}, \ref{def:gNW}, \ref{def:gAI} and \ref{def:gexp} with \emph{source} in place of \emph{edge}: $O_v^{\,i}$ is indexed by $i\in[t]^{E(v)}$, the symmetry group is $\prod_{e\in E}S_t$, and a set of copied observations is injectable if erasing indices is injective on it and any two of its members at vertices of a common source $e$ carry the same index on $e$. As in Definition~\ref{def:gAI}, the injectable sets are exactly the sets $\{O_v^{\,\iota|_{E(v)}}:v\in W\}$ with $W\subseteq V$ and $\iota:E(W)\to[t]$. We write $\mathcal J$ for any of $\INW$, $\IAI$, $\Iexp$.

\begin{lemma}[Facts that carry over]\label{lem:hcarry}
Let $H$ be a hypergraph scenario and $t\ge1$.
\begin{enumerate}[label=(\roman*),nosep]
\item $\Cscen{H}\subseteq\Iexp_t(H)\subseteq\IAI_t(H)\subseteq\INW_t(H)$.
\item $\Iexp_t(H)=\IAI_t(H)$.
\item If $t\ge2$, $P\in\INW_t(H)$ and $U,W\subseteq V$ satisfy $E(U)\cap E(W)=\varnothing$, then $P_{U\cup W}=P_U\otimes P_W$.
\item If $t\le T$ then $\mathcal J_T(H)\subseteq\mathcal J_t(H)$.
\end{enumerate}
\end{lemma}

\begin{proof}
(i) The proof of Proposition~\ref{prop:gsound} uses only that the sources are independent and that a copied observation reads one copy of each of its sources. (ii) The proofs of Lemmas~\ref{lem:trail} and~\ref{lem:rootsink} use only that every edge of the inflated causal graph joins a source copy, which has no parents, to a copied observation, which has no children. Here $S_e^{\,m}$ is a parent of $O_v^{\,i}$ exactly when $v\in e$ and $i_e=m$, so both proofs apply unchanged. (iii) A vertex in $U\cap W$ would have its sources in $E(U)\cap E(W)$, and it has at least one, so $U\cap W=\varnothing$. Let $\sigma$ transpose the indices $1$ and $2$ on every source of $E(W)$. It fixes $O_u^{\,1}$ for $u\in U$ and sends $O_w^{\,1}$ to $O_w^{\,2}$ for $w\in W$, so invariance equates the law of row $1$ on $U\cup W$, which is $P_{U\cup W}$, with the law that reads $U$ from row $1$ and $W$ from row $2$, which is $P_U\otimes P_W$ because distinct rows are independent. (iv) Keep the copied observations all of whose indices lie in $[t]$. The kept table is invariant under $\prod_eS_t$, embedded as the permutations fixing $t+1,\dots,T$, its first $t$ rows have law $P^{\otimes t}$, and injectable sets and ancestrally independent families at order $t$ are such sets and families at order $T$.
\end{proof}

\begin{lemma}[Absorption]\label{lem:absorption}
Let $e'\ne e$ be labels with $e'\subseteq e$, and let $H^-$ be $H$ with the label $e'$ deleted. Then $\Cscen{H^-}=\Cscen{H}$ and $\mathcal J_t(H^-)=\mathcal J_t(H)$ for every $t$ and every $\mathcal J$.
\end{lemma}

\begin{proof}
$H^-$ still covers $V$ because $e'\subseteq e$. Write $E^-=E\setminus\{e'\}$ and $E^-(v)$ for the sources of $H^-$ at $v$.

\emph{Compatible sets.} A model for $H^-$ is a model for $H$ whose responses ignore $S_{e'}$. Conversely, given a model for $H$, replace $S_e$ by the pair $(S_e,S_{e'})$ and delete $S_{e'}$. Every member of $e'$ is a member of $e$, so every response still reads its inputs, and the new sources are independent.

\emph{$\mathcal J_t(H^-)\subseteq\mathcal J_t(H)$.} Let $\Gamma^-$ witness $P$ for $H^-$. For $v\in e'$ split $i\in[t]^{E(v)}$ as $(i^-,i_{e'})$ with $i^-\in[t]^{E^-(v)}$, and for $v\notin e'$ put $i^-=i$. Let $\Gamma$ be the law of the family $O_v^{\,i}:=O_v^{\,i^-}$, the right side drawn from $\Gamma^-$. An element of $\prod_{E}S_t$ acts on this family as its restriction to $E^-$ acts on $\Gamma^-$, since the component on $e'$ changes nothing; so $\Gamma$ is invariant. Row $r$ of $H$ is row $r$ of $H^-$, so the diagonal law is $P^{\otimes t}$. An injectable set of $H$ with data $W$, $\iota$ consists of the variables $O_v^{\,\iota|_{E^-(v)}}$, $v\in W$, which form the injectable set of $H^-$ with data $W$, $\iota|_{E^-(W)}$; its law is $P_W$. Let $S_1,\dots,S_k$ be pairwise ancestrally independent injectable sets of $H$ with data $W_a$, $\iota_a$, so that $\iota_a(f)\ne\iota_b(f)$ for $a\ne b$ and $f\in E(W_a)\cap E(W_b)$. The corresponding sets of $H^-$ read subsets of the same source copies, hence are pairwise ancestrally independent. They are also pairwise disjoint: a common variable would sit at some $v\in W_a\cap W_b$ with $\iota_a$ and $\iota_b$ equal on $E^-(v)$, which is nonempty and on every member of which they differ. So their union has the product law under $\Gamma^-$, and $\Gamma$ satisfies every prescription of Definition~\ref{def:gAI} that $\Gamma^-$ satisfies. Lemma~\ref{lem:hcarry}(ii) gives the expressible case.

\emph{$\mathcal J_t(H)\subseteq\mathcal J_t(H^-)$.} Let $\Gamma$ witness $P$ for $H$. For $v\in e'$ and $j\in[t]^{E^-(v)}$ put $\widetilde O_v^{\,j}:=O_v^{\,(j,\,j_e)}$, which is defined because $v\in e$; for $v\notin e'$ put $\widetilde O_v^{\,j}:=O_v^{\,j}$. Distinct $j$ give distinct variables, so the law $\widetilde\Gamma$ of this family is a marginal of $\Gamma$. An element $\sigma^-$ of $\prod_{E^-}S_t$ extends to $\sigma\in\prod_ES_t$ by $\sigma_{e'}:=\sigma^-_e$, and $\sigma$ sends $O_v^{\,(j,j_e)}$ to $O_v^{\,(\sigma^-(j),\sigma^-(j)_e)}=\widetilde O_v^{\,\sigma^-(j)}$, so $\widetilde\Gamma$ is invariant. Rows go to rows. An injectable set of $H^-$ with data $W$, $\bar\iota$ consists of the same variables as the injectable set of $H$ with data $W$ and $\iota=\bar\iota$ extended by $\iota(e')=\bar\iota(e)$ when $e'\in E(W)$, which makes sense because a vertex of $W\cap e'$ lies in $e$. If $\bar\iota_a$ and $\bar\iota_b$ differ on $E^-(W_a)\cap E^-(W_b)$, the extensions also differ at $e'$ whenever $e'\in E(W_a)\cap E(W_b)$, since then $e$ lies in both. So every prescription for $H^-$ is a prescription for $H$.
\end{proof}

Deleting, one at a time, labels whose vertex set is contained in that of another label reaches a \emph{reduced} scenario, in which the vertex sets of distinct labels are distinct and pairwise incomparable, without changing the compatible set or any feasible set.

\begin{definition}[Trace]\label{def:trace}
For nonempty $U\subseteq V$ the \emph{trace} $H[U]$ has observers $U$ and the labels $e\in E$ with $e\cap U\ne\varnothing$, the label $e$ carrying the vertex set $e\cap U$.
\end{definition}

$H[U]$ covers $U$, and for $v\in U$ the sources of $H[U]$ at $v$ are the labels of $E(v)$. The inflations of $H$ and $H[U]$ therefore have the same index sets at every vertex of $U$, and a set of copied observations at vertices of $U$ is injectable in $H[U]$ if and only if it is injectable in $H$.

\begin{lemma}[Trace transport]\label{lem:htrace}
Let $\varnothing\ne U\subseteq V$, $t\ge1$, and fix a law $\nu_w$ on $\mathcal O_w$ for each $w\notin U$.
\begin{enumerate}[label=(\roman*),nosep]
\item If $Q\in\Cscen{H}$ then $Q_U\in\Cscen{H[U]}$.
\item If $P\in\mathcal J_t(H)$ then $P_U\in\mathcal J_t(H[U])$.
\item If $P_U\in\mathcal J_t(H[U])$ then $P:=P_U\otimes\bigotimes_{w\notin U}\nu_w\in\mathcal J_t(H)$.
\item With $P$ as in (iii), $d_{\TV}(P,\Cscen{H})\ge d_{\TV}(P_U,\Cscen{H[U]})$.
\end{enumerate}
\end{lemma}

\begin{proof}
(i) Keep the observations at $U$ of a model for $H$ and drop the sources that no retained observer reads. Each $v\in U$ reads the labels $E(v)$, which are its sources in $H[U]$.

(ii) Let $\Gamma$ witness $P$ for $H$ and let $\Gamma_U$ be its marginal on the copied observations at $U$. The group of $H[U]$ embeds in that of $H$ as the elements acting trivially on sources disjoint from $U$, with the same action on these variables, so $\Gamma_U$ is invariant. The rows restrict to rows, so the diagonal law is $P_U^{\otimes t}$. Injectable sets of $H[U]$ are injectable sets of $H$ with the same image, and source copies of $H[U]$ are source copies of $H$, so every prescription for $H[U]$ is a prescription for $H$.

(iii) Let $\Gamma^U$ witness $P_U$ for $H[U]$. Draw the copied observations at $U$ from $\Gamma^U$ and, independently, draw $O_w^{\,k}\sim\nu_w$ independently for all $w\notin U$ and $k\in[t]^{E(w)}$. An element of $\prod_ES_t$ acts on the first family through its components on the labels meeting $U$ and permutes the independent identically distributed variables at each $w\notin U$, so the table is invariant. Row $r$ at $U$ is row $r$ of $H[U]$ and the entries of the rows outside $U$ are distinct independent variables, so the diagonal law is $P^{\otimes t}$. Let $S$ be injectable in $H$ with data $W$, $\iota$. Its members at $U$ form the injectable set of $H[U]$ with data $W\cap U$, $\iota|_{E(W\cap U)}$, with law $(P_U)_{W\cap U}$; its other members are distinct independent variables with laws $\nu_w$; so its law is $P_W$. For pairwise ancestrally independent $S_1,\dots,S_k$, the parts at $U$ are pairwise ancestrally independent in $H[U]$, and two members of different $S_a$ at a common vertex $v$ carry different indices on every source of $E(v)\ne\varnothing$, so the parts at $U$ are pairwise disjoint and the variables outside $U$ are distinct. The joint law is therefore the product of the laws $P_{W_a}$. Lemma~\ref{lem:hcarry}(ii) gives the expressible case.

(iv) For $Q\in\Cscen{H}$, part (i) gives $Q_U\in\Cscen{H[U]}$, and $d_{\TV}(P,Q)\ge d_{\TV}(P_U,Q_U)\ge d_{\TV}(P_U,\Cscen{H[U]})$.
\end{proof}

\begin{corollary}[Traces and blocks]\label{cor:hblocks}
\begin{enumerate}[label=(\roman*),nosep]
\item If $\mathcal J_T(H)=\Cscen{H}$ then $\mathcal J_T(H[U])=\Cscen{H[U]}$ for every nonempty $U\subseteq V$.
\item Let $V=V_1\sqcup\dots\sqcup V_k$ with no source meeting two parts, and suppose $\mathcal J_{T_a}(H[V_a])=\Cscen{H[V_a]}$ for each $a$. Then $\mathcal J_T(H)=\Cscen{H}$ for $T=\max\{2,T_1,\dots,T_k\}$.
\end{enumerate}
\end{corollary}

\begin{proof}
(i) If $P_U\in\mathcal J_T(H[U])$, Lemma~\ref{lem:htrace}(iii) puts $P_U\otimes\bigotimes_w\nu_w$ in $\mathcal J_T(H)=\Cscen{H}$, and Lemma~\ref{lem:htrace}(i) puts $P_U$ in $\Cscen{H[U]}$. (ii) Let $P\in\mathcal J_T(H)$. The sets $E(V_a)$ are pairwise disjoint and $T\ge2$, so Lemma~\ref{lem:hcarry}(iii), iterated, gives $P=\bigotimes_aP_{V_a}$. Lemma~\ref{lem:htrace}(ii) and Lemma~\ref{lem:hcarry}(iv) give $P_{V_a}\in\mathcal J_{T_a}(H[V_a])=\Cscen{H[V_a]}$. Models for the blocks on disjoint source families form a model for $H$, since $H[V_a]$ consists of the sources of $H$ inside $V_a$, and its law is $P$. With Lemma~\ref{lem:hcarry}(i), $\mathcal J_T(H)=\Cscen{H}$.
\end{proof}

\subsection{Termination}

An observer who reads every source does not make every law compatible: in the star with sources $\{b,a_1\}$ and $\{b,a_2\}$ the observer $b$ reads both sources, yet $O_{a_1}$ and $O_{a_2}$ are independent in every model. What such an observer does is drop out of the problem.

\begin{theorem}[Observers that read every source]\label{thm:hdominating}
Let $b\in V$ lie in every source of $H$, and let $V\ne\{b\}$. Then
\[
\Cscen{H}=\bigl\{P:\ P_{V\setminus\{b\}}\in\Cscen{H[V\setminus\{b\}]}\bigr\},
\]
and for every $t$ and every $\mathcal J$, $\mathcal J_t(H)=\Cscen{H}$ if and only if $\mathcal J_t(H[V\setminus\{b\}])=\Cscen{H[V\setminus\{b\}]}$.
\end{theorem}

\begin{proof}
Write $H'=H[V\setminus\{b\}]$. The inclusion $\subseteq$ is Lemma~\ref{lem:htrace}(i). For $\supseteq$, let $P_{V\setminus\{b\}}\in\Cscen{H'}$ and take a model for it with deterministic responses. Run it on the sources of $H$, giving a trivial value to a source whose vertex set is $\{b\}$ if there is one; each $v\ne b$ reads the same labels as in $H'$. Append to one source an independent seed $\zeta$ uniform on $[0,1]$, ignored by every $v\ne b$. The observer $b$ reads every source, so $(O_v)_{v\ne b}$ and $\zeta$ are functions of what $b$ reads, and $O_b:=\varphi\bigl((O_v)_{v\ne b},\zeta\bigr)$, with $\varphi(o,\cdot)$ pushing the uniform law to $P(O_b\in\cdot\mid O_{V\setminus\{b\}}=o)$, is a legitimate response. The observed law is $P$.

If $\mathcal J_t(H')=\Cscen{H'}$ and $P\in\mathcal J_t(H)$, Lemma~\ref{lem:htrace}(ii) gives $P_{V\setminus\{b\}}\in\Cscen{H'}$, hence $P\in\Cscen{H}$. The converse is Corollary~\ref{cor:hblocks}(i).
\end{proof}

For three observers reading one common source, Theorem~\ref{thm:hdominating} applied twice reduces the scenario to one observer with its own source, where every law is compatible and the order-one test is exact; this is Remark~\ref{rem:hypergraph}.

\begin{definition}[Hyper-double-star]\label{def:hds}
Let $H$ be connected, that is, the relation \emph{some source contains both} generates a single class on $V$. Its \emph{centres} are $Z=\{v:\deg v\ge2\}$ and its \emph{leaves} are the other observers. $H$ is a \emph{hyper-double-star} if $Z\subseteq y$ for some source $y$. A source $e$ is \emph{inner} if $Z\subseteq e$ and \emph{split} otherwise; write $L_e=e\setminus Z$ for its leaves.
\end{definition}

Every leaf lies in exactly one source, so the sets $L_e$ are pairwise disjoint and cover the leaves; $L_e$ may be empty.

\begin{theorem}[Hyper-double-star reconstruction]\label{thm:hds}
Let $H$ be a connected hyper-double-star with finite observed alphabets, and put
\[
T_H=\max\Bigl(\{2\}\cup\Bigl\{\prod_{l\in L_e}|\mathcal O_l|:\ e\text{ split}\Bigr\}\Bigr).
\]
Then $\mathcal J_T(H)=\Cscen{H}$ for every $T\ge T_H$ and every $\mathcal J$.
\end{theorem}

\begin{proof}
By Lemma~\ref{lem:hcarry}(i) it suffices to show $\INW_T(H)\subseteq\Cscen{H}$ for $T\ge T_H$. Let $\Gamma$ witness $P\in\INW_T(H)$. Write $I$ and $S$ for the sets of inner and split sources; $y\in I$. For a family $\xi=(\xi_e)_e$ of leaf values write $\{O_L=\xi\}$ for the event that $O_{L_e}=\xi_e$ for every $e$ in the index set of $\xi$.

\emph{(i) Leaves.} If $L_e\ne\varnothing$ then $E(L_e)=\{e\}$, so the sets $E(L_e)$ are pairwise disjoint and Lemma~\ref{lem:hcarry}(iii), iterated, gives that $P_L$, the law of all leaves, is $\bigotimes_eP_{L_e}$. A leaf $l\in L_e$ has the copies $O_l^{\,m}$, $m\in[T]$, and $O_l^{\,m}$ belongs to row $m$. Since the rows are independent with law $P$, the blocks $(O_l^{\,m})_{l\in L_e}$ are independent over all pairs $(e,m)$, with laws $P_{L_e}$.

\emph{(ii) Conditioning events.} For $e\in S$ list $x_e^{\,1},\dots,x_e^{\,T}\in\prod_{l\in L_e}\mathcal O_l$ so that every element of $\supp P_{L_e}$ occurs, which $T\ge T_H$ allows, and fill the remaining places with supported values. Fix $p_e:\supp P_{L_e}\to[T]$ with $x_e^{\,p_e(\xi_e)}=\xi_e$. For $u=(u_e)_{e\in I}\in\prod_{e\in I}\supp P_{L_e}$ let
\[
E_u=\bigcap_{e\in I,\,m\in[T]}\{O_{L_e}^{\,m}=u_e\}\ \cap\bigcap_{e\in S,\,m\in[T]}\{O_{L_e}^{\,m}=x_e^{\,m}\}.
\]
By (i), $\Gamma(E_u)=\prod_{e\in I}P_{L_e}(u_e)^T\prod_{e\in S}\prod_mP_{L_e}(x_e^{\,m})>0$.

\emph{(iii) The finite-order identity.} Fix $n\in[T]^S$ and let $\iota_n:E\to[T]$ equal $n_e$ on $S$ and $1$ on $I$. Choose $\sigma_e\in S_T$ with $\sigma_e(1)=n_e$ for $e\in S$ and $\sigma_e$ the identity for $e\in I$. The resulting symmetry carries row $m$ to the family with centre entries $O_z^{\,(\sigma_e(m))_{e\in E(z)}}$ and leaf entries $O_l^{\,\sigma_e(m)}$, $l\in L_e$; for $m=1$ the centre entries are $O_z^{\,\iota_n|_{E(z)}}$. By invariance the $T$ images have joint law $P^{\otimes T}$. Every copied leaf observation occurs in exactly one image, and $E_u$ says that in image $m$ the leaves of $L_e$ take the value $u_e$ for $e\in I$ and $x_e^{\,\sigma_e(m)}$ for $e\in S$. Marginalizing images $2,\dots,T$ over their centre entries and using (i),
\[
\Gamma\bigl(O_Z^{\,\iota_n}=\beta,\ E_u\bigr)=P\bigl(O_Z=\beta,\ O_{L}=(u,x^n)\bigr)\prod_{e\in I}P_{L_e}(u_e)^{T-1}\prod_{e\in S}\prod_{m\ne n_e}P_{L_e}(x_e^{\,m}),
\]
where $O_Z^{\,\iota_n}=(O_z^{\,\iota_n|_{E(z)}})_{z\in Z}$ and $x^n=(x_e^{\,n_e})_{e\in S}$. Dividing by $\Gamma(E_u)$ and using $P_L=\bigotimes_eP_{L_e}$,
\begin{equation}\label{eq:hcond}
\Gamma\bigl(O_Z^{\,\iota_n}=\beta\bigm|E_u\bigr)=P\bigl(O_Z=\beta\bigm|O_L=(u,x^n)\bigr).
\end{equation}

\emph{(iv) Response tables.} For $z\in Z$ let $S(z)=S\cap E(z)$; every source at $z$ lies in $I$ or in $S(z)$. For $\xi\in\prod_{e\in S(z)}\supp P_{L_e}$ let $F_z(\xi)$ be the copied observation $O_z^{\,j}$ with $j_e=p_e(\xi_e)$ for $e\in S(z)$ and $j_e=1$ for $e\in I$. It is indexed by the values on the split sources that $z$ reads and by nothing else, because the positions $p_e$ are chosen one source at a time. If $n_e=p_e(\xi_e)$ for all $e\in S$ then $O_z^{\,\iota_n|_{E(z)}}=F_z(\xi|_{S(z)})$ for every $z$. Let $\mu_u$ be the law of the family $(F_z(\xi))_{z,\xi}$ under $\Gamma(\cdot\mid E_u)$, and let $\Theta^u=(\Theta^u_z(\xi))_{z,\xi}$ denote a random family with law $\mu_u$. By \eqref{eq:hcond}, for every $\xi\in\prod_{e\in S}\supp P_{L_e}$,
\begin{equation}\label{eq:hmu}
\mu_u\bigl(\Theta^u_z(\xi|_{S(z)})=\beta_z\ \forall z\in Z\bigr)=P\bigl(O_Z=\beta\bigm|O_L=(u,\xi)\bigr).
\end{equation}
Now take independent sources: $X_e\sim P_{L_e}$ for $e\ne y$, and at $y$ the pair $\bigl(X_y,(\Theta^u)_u\bigr)$ with $X_y\sim P_{L_y}$ and $\Theta^u\sim\mu_u$, all independent. Respond by $O_l=(X_e)_l$ for $l\in L_e$ and
\[
O_z=\Theta^{u}_z\bigl((X_e)_{e\in S(z)}\bigr),\qquad u=(X_e)_{e\in I},
\]
for $z\in Z$. A leaf reads its own source. A centre reads every inner source, $y$ among them, and the split sources $S(z)$, which is everything its response uses. The leaves have law $\bigotimes_eP_{L_e}=P_L$, and given $O_L=(u,\xi)$ the centres have the law \eqref{eq:hmu}, because $\Theta^u$ is independent of the $X_e$. Leaf values outside the supports have probability zero under both laws. So the observed law is $P$.
\end{proof}

A leaf-free source ($L_e=\varnothing$) contributes the factor $1$ and the single empty value. The proof does not use that $H$ is reduced. The positions $p_e$ must be chosen source by source: a choice of $n$ depending jointly on all leaf values would give a table $F_z$ that reads leaves of sources $z$ does not read. Inner sources need no listing, since conditioning all their leaf copies on one value is enough; this is why a star needs only order two, whatever its alphabets.

\begin{corollary}[Peeling]\label{cor:hpeel}
Let $\mathcal P$ be the smallest class of scenarios, closed under isomorphism, such that $H\in\mathcal P$ whenever each connected component $K$ of the reduction of $H$ is a hyper-double-star or has a nonempty set $D_K$ of observers lying in every source of $K$ with $K[V(K)\setminus D_K]\in\mathcal P$. If $H\in\mathcal P$ then $\mathcal J_T(H)=\Cscen{H}$ for some finite $T$, namely $T=\max\{2,T_{K_1},\dots,T_{K_r}\}$, where $K_1,\dots,K_r$ are the hyper-double-stars at which the recursion stops.
\end{corollary}

\begin{proof}
Induction on $|V|$. Absorption (Lemma~\ref{lem:absorption}) replaces $H$ by its reduction. A component $K$ that is a hyper-double-star has $\mathcal J_{T_K}(K)=\Cscen{K}$ by Theorem~\ref{thm:hds}. Otherwise $K$ has at least two sources, and $V(K)\setminus D_K\ne\varnothing$ since in a reduced scenario no source equals $D_K$ unless it is the only one. Each $b\in D_K$ lies in every source of each successive trace, so Theorem~\ref{thm:hdominating}, applied once for each element of $D_K$, transfers termination at any order from $K[V(K)\setminus D_K]$, where it holds by induction, to $K$. Corollary~\ref{cor:hblocks}(ii) assembles the components.
\end{proof}

\begin{proposition}[Graphs]\label{prop:hgraph}
Let $G$ be a connected simple graph without isolated vertices, read as a pair-source scenario. Then $G$ is a hyper-double-star if and only if it is a double-star, and $G\in\mathcal P$ if and only if it is a double-star. If $G$ is a double-star with central edge $\{b,c\}$ then $T_G\le\max(\{2\}\cup\{|\mathcal O_v|:v\ne b,c\})$, the order of Theorem~\ref{thm:doublestar}, with equality when $b$ and $c$ both have degree at least two.
\end{proposition}

\begin{proof}
Here $Z$ is the set of vertices of degree at least two and sources are edges, so $G$ is a hyper-double-star exactly when $|Z|\le2$ and the vertices of $Z$ are adjacent when there are two. If $Z=\varnothing$ the connected graph $G$ is one edge. If $Z=\{b\}$ every other vertex has degree one, and it is adjacent to $b$, since two adjacent degree-one vertices would form a component; so $G$ is a star. If $Z=\{b,c\}$ with $bc$ an edge, every other vertex has degree one and is adjacent to $b$ or $c$ by the same argument, so $G$ is a double-star. Conversely every vertex of a double-star other than $b,c$ has degree one. For $\mathcal P$: a vertex lying in every edge of a connected graph makes it a star, which is already a hyper-double-star. For the order: if $b$ and $c$ both have degree at least two, $Z=\{b,c\}$, the only inner source is $bc$, and the split sources are the leaf edges, each with a single leaf, so $T_G$ is the displayed maximum. Otherwise $|Z|\le1$, every edge is inner and $T_G=2$.
\end{proof}

So Theorem~\ref{thm:hds} contains Theorem~\ref{thm:doublestar}, with the same order when both centres have degree at least two and with order two for stars, in agreement with \cite[Section 4.1]{NW2020}.

\subsection{Nontermination through pair sets}

\begin{definition}[Pair set]\label{def:pairset}
$U\subseteq V$ is a \emph{pair set} of $H$ if $|e\cap U|\le2$ for every source $e$. Its graph $G_U$ has vertex set $U$ and an edge $e\cap U$ for every source with $|e\cap U|=2$, repetitions discarded.
\end{definition}

\begin{theorem}[Nontermination transported through pair sets]\label{thm:hpairset}
Let $H$ have binary observations, and suppose some pair set $U$ has a connected component of $G_U$ with at least one edge that is not a double-star. Then for every $t\ge1$ there is a strictly positive rational law
\[
P\in\Iexp_t(H)\setminus\Cscen{H}\subseteq\IAI_t(H)\setminus\Cscen{H}\subseteq\INW_t(H)\setminus\Cscen{H},
\]
with $d_{\TV}(P,\Cscen{H})\ge11/(640\,m^2t^2)$ if that component has an induced cycle of length $m$, and $d_{\TV}(P,\Cscen{H})\ge1/(1536\,t^2)$ if it has an induced path on five vertices. In particular no finite order of any of the three hierarchies characterizes $\Cscen{H}$.
\end{theorem}

\begin{proof}
Let $K$ be such a component. By Lemma~\ref{lem:exhaustion} it has an induced subgraph $R$ isomorphic to $C_m$ with $m\ge3$ or to $P_5$; take $R$ to be the induced subgraph named in the distance bound being proved, and let $U'=V(R)$. Then $U'\subseteq U$ is a pair set and $G_{U'}=R$: a source with $|e\cap U'|=2$ has $e\cap U=e\cap U'$ because $|e\cap U|\le2$, so $e\cap U'$ is an edge of $G_U$ with both ends in $U'$, hence an edge of $K$ and of the induced subgraph $R$; conversely each edge of $R$ is such a trace. In $H[U']$ every source has one or two members, those with two members are the edges of $R$, possibly repeated, and every vertex of $U'$ lies on an edge of $R$. Lemma~\ref{lem:absorption} deletes the repetitions and the one-member sources, so $\Cscen{H[U']}=\Cscen{R}$ and $\mathcal J_t(H[U'])=\mathcal J_t(R)$. Theorem~\ref{thm:cycle} or Theorem~\ref{thm:fivepath} gives a strictly positive rational $P_R\in\Iexp_t(R)\setminus\Cscen{R}$ with the displayed distance from $\Cscen{R}$. Put $P=P_R\otimes\Bern(1/2)^{\otimes(V\setminus U')}$. Lemma~\ref{lem:htrace}(iii) gives $P\in\Iexp_t(H)$, Lemma~\ref{lem:htrace}(iv) gives the distance, and Lemma~\ref{lem:hcarry}(i) gives the other two inclusions.
\end{proof}

By Lemma~\ref{lem:exhaustion} the hypothesis holds if and only if some pair set $U$ has $G_U$ isomorphic to a cycle or to $P_5$. It can hold for no pair set at all. If every three observers lie together in some source, every pair set has at most two elements and Theorem~\ref{thm:hpairset} says nothing.

\subsection{A candidate classification}

\begin{conjecture}\label{conj:hypergraph}
A hypergraph scenario with binary observations admits a finite characterizing order, for any and then for all of the three hierarchies, if and only if it lies in the class $\mathcal P$ of Corollary~\ref{cor:hpeel}.
\end{conjecture}

The \emph{if} direction is Corollary~\ref{cor:hpeel}. The \emph{only if} direction holds whenever some pair set carries a non-double-star component (Theorem~\ref{thm:hpairset}); the two results never apply to the same scenario, as they must not, and an exhaustive enumeration of the reduced scenarios on at most five observers, in \texttt{artifact/certificates/hypergraph/} of the repository \repo, confirms this. Proposition~\ref{prop:hgraph} and Theorem~\ref{thm:classification} show that on pair-source scenarios the conjecture is a theorem.

The recursion through Theorem~\ref{thm:hdominating} cannot be dropped. The scenario with observers $b,1,\dots,6$ and sources $\{b,1,2\},\{b,2,3\},\{b,4,5\},\{b,5,6\}$ has centres $\{b,2,5\}$, which no source contains, so it is not a hyper-double-star. Deleting $b$ leaves two paths $P_3$, so by Theorem~\ref{thm:hdominating} and Theorem~\ref{thm:classification} it terminates at order two. Seven observers is the least number for which a connected scenario lies in $\mathcal P$ without being a hyper-double-star. Let $K$ be such a scenario, reduced, and $D=D_K$. Deleting $D$ leaves the degrees of the other observers unchanged, and every observer of $D$ is a centre. If every component of $K[V(K)\setminus D]$ were a hyper-double-star and at most one had a centre, its centres would lie in a source $y\setminus D$ with $y$ a source of $K$, and then the centres of $K$ would lie in $y$. So either some component is itself in $\mathcal P$ without being a hyper-double-star, and has at least seven observers by induction, or two components have centres. A component with a centre contains two sources sharing an observer, neither containing the other, hence at least three observers; with $D\ne\varnothing$ this gives at least seven.

The smallest scenario the two results leave open is the complete $3$-uniform hypergraph on four observers,
\[
K_4^{(3)}:\qquad E=\bigl\{\{1,2,3\},\{1,2,4\},\{1,3,4\},\{2,3,4\}\bigr\}.
\]
No observer reads every source, and the centres $\{1,2,3,4\}$ lie in no source, so $K_4^{(3)}\notin\mathcal P$. Its pair sets have at most two elements. Its trace on any three observers contains a source equal to all three, so every law on three of the observers is compatible, and no restriction to fewer observers can witness nontermination. An exhaustive enumeration of the twenty reduced scenarios on four observers, up to isomorphism, finds that every one other than $K_4^{(3)}$ is decided by Corollary~\ref{cor:hpeel} or Theorem~\ref{thm:hpairset}. Whether $K_4^{(3)}$ terminates is open.
